% TOP_PROBLEM: {"id": 3163, "title": "(m,n)-cycle covers", "category_id": 3, "category": {"id": 3, "name": "graph_theory", "display_name": "Graph Theory", "description": "Problems involving graphs, networks, and their properties.", "slug": "graph-theory", "order_index": 3, "created_at": "2026-07-31T15:26:25.671Z"}, "status": "open", "problem_number": "OPG-139", "set_id": 12}
% FIRST_POSTED: 2026-10-01
% ATTEMPT_STATUS: unresolved
% UnsolvedMath ID 3163; source catalogue code OPG-139
% !TeX program = lualatex
% GPT-6 Astra Ultra research attempt; independent partial work, 2026-10-01.
\documentclass[11pt,a4paper]{article}
\usepackage[margin=25mm]{geometry}
\usepackage{fontspec}
\setmainfont{lmroman10-regular.otf}[BoldFont=lmroman10-bold.otf,
 ItalicFont=lmroman10-italic.otf,BoldItalicFont=lmroman10-bolditalic.otf]
\usepackage{amsmath,amssymb,mathtools}
\usepackage{unicode-math}
\setmathfont{latinmodern-math.otf}
\AtBeginDocument{\let\setminus\smallsetminus}
\usepackage{xurl,hyperref}
\hypersetup{colorlinks=true,urlcolor=blue}
\setcounter{secnumdepth}{0}
\setlength{\parindent}{0pt}
\setlength{\parskip}{5pt}
\setlength{\emergencystretch}{3em}
\begin{document}
\section*{(m,n)-cycle covers}\label{3163}
{\small UnsolvedMath ID 3163 · OPG-139 · Graph Theory · OpenGarden}
\subsection{Definitions and mathematical statement}
Let $G=(V,E)$ be a finite undirected loopless multigraph.
A binary cycle is an edge set $C\subseteq E$ such that every
vertex of the spanning subgraph $(V,C)$ has even degree.
The empty set is allowed. A binary cycle may be disconnected
and may have vertices of degree greater than two; equivalently,
its edges can be partitioned into ordinary circuits.

For positive integers $m,n$, an $(m,n)$-cycle cover is a list
$C_1,\ldots,C_m$ of binary cycles such that
\[
 \sum_{i=1}^m\mathbf1_{\{e\in C_i\}}=n
 \qquad\text{for every }e\in E.
\]
The list may contain repeated entries, and empty entries allow
a cover with fewer than $m$ nonempty members to be padded.
Membership of one edge in one binary cycle counts once.

The conjecture is that every bridgeless $G$ admits a
$(5,2)$-cycle cover, where bridgeless means that deleting any
single edge does not increase the number of connected
components. Thus five even-degree edge sets should together
cover each original edge exactly twice.

The number five bounds these binary cycles, not the number
of connected circuits arising after they are decomposed.
In a cubic graph every nonzero degree inside a binary cycle
is two, so each of the five members is a vertex-disjoint union
of circuits. No connectivity of $G$ or of the covering members
is required; components can be handled with the same five
indices.
\subsection{Short English statement}
Can the edges of every bridgeless graph be covered exactly twice by five even-degree subgraphs?
\subsection{Research attempt: compressing four binary cycles and the five-colour barrier}
\paragraph{Scope and literature check.}
GPT-6 Astra Ultra, 1 October 2026. Here a cycle means an even-degree
edge set, not necessarily one circuit. Oum's 2026 exposition [S3]
presents the announced ordinary cycle double cover proof and explicitly
lists the five-cycle double cover conjecture separately. Thus a cover
with an unrestricted number of circuits does not already settle the
canonical five-member target. Liu, Hao, Luo and Zhang [S4] study
five-cycle covers using flows and reductions. The work below supplies
a self-contained compression argument, identifies its exact limit, and
checks a small obstruction to using fewer members.

\paragraph{1. Three members and a two-coordinate flow.}
Suppose $A,B,C$ are binary cycles covering each edge twice. The
membership sum is zero modulo two, so $C=A\mathbin{\triangle}B$.
No edge is absent from both $A$ and $B$, since then it could occur
only in $C$, contrary to multiplicity two. Consequently the labels
\[
 f(e)=(\mathbf1_A(e),\mathbf1_B(e))\in\mathbb F_2^2
\]
are nonzero, and their sum at every vertex is zero because $A$ and
$B$ are even subgraphs. Conversely, any such nonzero labelling with
zero vertex sums gives the binary cycles $A,B,A\triangle B$, covering
each edge twice: a nonzero two-bit vector has total membership two
in its first coordinate, second coordinate, and their sum.

\paragraph{2. Four members can always be compressed to three.}
Let $C_1,\ldots,C_4$ be a double cover. Assign distinct labels
$t_1=00,t_2=01,t_3=10,t_4=11$ in $\mathbb F_2^2$. Each edge belongs
to exactly two distinct indices $i,j$, even if some members of the
list happen to be equal as sets. Label that edge $f(e)=t_i+t_j$.
It is nonzero. At each vertex $v$,
\[
 \sum_{e\ni v}f(e)=\sum_{i=1}^4 t_i\deg_{C_i}(v)=0,
\]
since each degree on the right is even. The preceding construction
therefore produces a three-member double cover. Padding with an empty
member proves the converse. Hence $(3,2)$ and $(4,2)$ existence are
equivalent for every loopless multigraph.

For cubic graphs the three nonzero labels on incident edges must be
pairwise distinct: if two were equal, the zero-sum condition would
force the third to be zero. Thus a three-member double cover gives
a proper three-edge-colouring. Conversely, the three unions of pairs
of colour classes form a double cover. This establishes the equivalence
without confusing connected circuits with binary cycles.

\paragraph{3. Why the same compression does not handle five.}
An assignment of fixed labels $t_1,\ldots,t_5\in\mathbb F_2^2$ cannot
make $t_i+t_j$ nonzero for all distinct pairs: there are only four
labels, so some pair coincides. A particular cover might avoid that
pair, but the five-member hypothesis alone gives no such restriction.
This disproves only the universal fixed-label compression scheme. It
does not disprove a five-member double cover or rule out a construction
using a different algebraic structure.

A concrete check is provided by the Petersen graph, with outer cycle
$0,1,2,3,4,0$, spokes $i(i+5)$, and inner edges
$(i+5)(5+(i+2\bmod5))$. Exhaustively checking its $\binom{15}{5}$
five-edge subsets finds exactly six perfect matchings; every complement
consists of two five-cycles. A proper three-edge-colouring would make
each colour class a perfect matching whose complement is a union of
even alternating circuits. Hence this graph has no three- or
four-member double cover. The check script [S5] reproduces that entire
finite argument, including the component lengths, and verifies an
explicit five-member double cover. Thus the obstruction does not extend
to the conjectured number five.

\paragraph{Verification and first remaining gap.}
All parity identities above hold for parallel edges counted separately.
The finite Petersen enumeration certifies one graph only; it is not
evidence for a uniform bound over arbitrary order. The unresolved task
is producing five even edge sets on every bridgeless graph, or an
appropriate pair-label system satisfying every vertex constraint. Neither
the ordinary double-cover theorem nor the two-coordinate compression
provides that construction.

\subsection{Final status}
UNRESOLVED. The equivalence of three- and four-member covers is proved,
the fixed-label five-member shortcut is obstructed, and the Petersen
boundary example is verified. The general five-member assertion is
not established. Completion estimate: no defensible global percentage.

\subsection{Sources}
\begin{itemize}
\item[S1] ULAM.AI and the UnsolvedMath Contributors, supplied problems.json, record 3163 (OPG-139), OpenGarden collection. Catalogue identity and source transcription; CC BY 4.0.
\url{https://huggingface.co/datasets/ulamai/UnsolvedMath}.
\item[S2] Open Problem Garden, (m,n)-cycle covers. Original problem and discussion; the mathematical formulation below is expanded from the supplied source record.
\url{http://www.openproblemgarden.org/op/m_n_cycle_covers}.
\item[S3] S.-i. Oum, \emph{A proof of the cycle double cover conjecture by OpenAI: An exposition}, arXiv:2607.16356v3 (2026), Section 9.2. Checked 1 October 2026.
\url{https://arxiv.org/abs/2607.16356}.
\item[S4] S. Liu, R.-X. Hao, R. Luo and C.-Q. Zhang, \emph{5-Cycle Double Covers, 4-Flows, and Catlin Reduction}, SIAM Journal on Discrete Mathematics (2023). Checked 1 October 2026.
\url{https://doi.org/10.1137/22M1472425}.
\item[S5] GPT-6 Astra Ultra, exhaustive Petersen matching check and five-cover certificate, 1 October 2026: \texttt{scripts/check\_attempt\_batch02\_connectivity\_2026\_10\_01.py} in this repository; run with Python 3.
\end{itemize}
\end{document}
